\section{Introduction}
\label{sec:intro}

\IEEEPARstart{U}{nmanned} aerial vehicles (UAVs) equipped with gas sensors have emerged as promising platforms for autonomous gas source localization (GSL), with applications in industrial leak detection, environmental monitoring, and disaster response~\cite{ref1,ref2}. The fundamental challenge is that gas dispersion in complex environments is driven by turbulent transport, resulting in intermittent, filamentary plume structures~\cite{ref3,ref4}. A sensor may read high concentrations at one moment and near-zero readings seconds later, even without the UAV moving.

\subsection{Limitations of Existing Methods}

Current GSL methods fall into three families, each with distinct failure modes:

\textbf{Upwind-tracking methods} (e.g., SurgeCast~\cite{ref5}) track the gas plume upwind toward the source. These achieve excellent performance when the UAV starts within a strong plume (0.4m mean error), but degrade dramatically when the initial signal is weak. In our experiments, SurgeCast's error increases by $8\times$ (from 0.4m to 3.3m) when the starting position has weak plume contact---a common scenario where the UAV's initial position is arbitrary.

\textbf{Plume-model methods} (e.g., PMFS~\cite{ref6}, GrGSL~\cite{ref7}) fit parametric dispersion models to observed concentrations. When the actual dispersion deviates from model assumptions---as with filament-based transport---these methods systematically fail (3--8m error across all conditions).

\textbf{Information-theoretic methods} (e.g., Infotaxis~\cite{ref8}) select actions to maximize expected information gain. While theoretically elegant, standard Infotaxis lacks a principled stopping criterion: the algorithm has no mechanism to declare ``the source has been found'' with any confidence guarantee.

\subsection{The Stopping Problem}

A critical and under-studied challenge is deciding \emph{when to stop searching}. Existing methods either run for a fixed number of steps or use ad-hoc convergence criteria (e.g., ``MAP hasn't moved for $K$ steps'') without formal guarantees. In safety-critical applications, premature termination leads to false source declarations; unnecessarily prolonged searches waste time and battery. Neither failure mode is acceptable, yet no existing GSL method provides formal guarantees on false termination probability.

Conformal prediction (CP)~\cite{ref9} has been proposed for distribution-free inference, but its exchangeability assumption is violated in sequential Bayesian inference where the posterior changes at each step. This fundamental mismatch means CP-based stopping rules cannot provide the guarantees they claim.

\subsection{Our Approach}

We address these challenges through a unified framework based on two recent advances:

\textbf{Anytime-Valid Inference (AVI)}~\cite{ref10,ref11} provides sequential hypothesis testing where the stopping time need not be pre-specified. At its core are \emph{e-processes}---products of non-negative random variables (e-values) that maintain validity under arbitrary stopping rules, without exchangeability assumptions. This makes AVI particularly suited for GSL.

\textbf{Blackwell Approachability}~\cite{ref12,ref13} provides a game-theoretic framework for characterizing achievable objectives in sequential decision problems. By modeling the UAV-environment interaction as a repeated vector-valued game, we jointly optimize search efficiency, stopping decisions, and risk sensitivity.

\subsection{Contributions}

\begin{enumerate}
\item \textbf{Unified Framework.} We formulate UAV GSL as a sequential decision problem and prove that the coverage-efficiency target set is approachable under mild conditions (Proposition~\ref{prop:approachability}).

\item \textbf{Anytime-Valid Stopping.} We construct an e-process from posterior convergence indicators and prove $\mathbb{P}(\text{false positive}) \leq \delta$ at any stopping time (Theorem~\ref{thm:anytime}), without exchangeability assumptions. This is, to our knowledge, the first application of anytime-valid inference to robotic source localization.

\item \textbf{Risk-Adjusted Search.} We employ risk-adjusted information-directed sampling and prove it accelerates e-process convergence (Proposition~\ref{prop:coupling}).

\item \textbf{Experimental Validation.} On three indoor environments with varying signal strengths, our method achieves consistent performance (1.2--2.1m) with $2\times$ lower cross-environment variance than SurgeCast, while the e-process fires in 100\% of runs.
\end{enumerate}